\begin{theorem}[the remaining scalar is nonzero]
\label{thm:h12-scalar}
\Verified\ The induced map
\[
 H^0(B,R^1f_*\mathcal E)\simeq\mathbb C^2
 \xrightarrow{\ H^0(\kappa_1)\ }
 H^0(B,\mathcal O)\simeq\mathbb C
\]
has rank one.
\end{theorem}

\begin{proof}
Every matrix $R_g$ in (2.3d) has second column $(0,1)^{\mathsf t}$.
Consequently the marked vertical field
\[
 \xi=\frac{\partial}{\partial\zeta_2}
\]
is invariant under every deck transformation.  The affine translations in
\cref{prop:candidate-finite-fillings} have identity differential in this
direction, and the corresponding toric logarithmic direction extends across
the cusp charts of \cref{thm:candidate-cusp-geometry}.  Thus $\xi$ is a
global section nowhere zero in the locally free relative sheaf
$T_{X/B}=\mathcal E^\vee$ and is the section used in the evaluation sequence
\[
 0\longrightarrow\mathcal O(-1)\longrightarrow f_*\mathcal E
 \xrightarrow{\operatorname{ev}_\xi}\mathcal O\longrightarrow0.
\]
This sequence splits, and the lift $h$ of $1\in H^0(\mathcal O)$ is unique:
the difference of two lifts is a global section of $\mathcal O(-1)$, while
$H^0(B,\mathcal O(-1))=0$.  On the marked cover write that unique lift as
\[
 h=(a(z),1)\binom{d\zeta_1}{d\zeta_2}.
\tag{11.18}
\]
The function $a$ is holomorphic on each finite cover because $h$ is a
section of the actual invariant lattice, not merely of its generic fibre.
For $g_2$, invariance of (11.18) and
$R_{g_2}=\left(\begin{smallmatrix}1/\tau&0\\-\mu/\tau&1\end{smallmatrix}\right)$
mean exactly
\[
 (a(z),1)=(a(g_2z),1)R_{g_2}(z),
 \qquad\text{hence}\qquad
 a(g_2z)=\tau(z)a(z)+\mu(z).
\tag{11.18a}
\]
At the fixed point $z_2$, the normalized values
$\tau(z_2)=i$ and $\mu(z_2)=(1-i)/2$ therefore give
\[
 a(z_2)=\frac12.
\tag{11.18b}
\]
The cup-product map
$f_*\mathcal E\otimes G\to R^1f_*\mathcal E$ sends
$h\otimes\gamma$ to the unit generator $e_0b_0$ of the degree-zero
character line in (11.4).  This identifies the precise entry of the row
(11.11) that must be evaluated.

Normalize the order-four linearizing coordinate by
\[
 s=(z_2-\bar z_2)s_2,
 \qquad s\circ g_2=-is,
 \qquad \frac{ds}{dz}(z_2)=1.
\]
Write a dot for $d/ds$, and put
$\ell=dt/t=4\,ds/s$.  In the quotient basis $(\gamma,u)$ of $G$, the
marked period derivative and $ds=(s/4)\ell$ give the exact logarithmic map
\[
 \kappa_0(h)=\frac{s}{4}
 \Bigl[(6a\dot\mu+\dot\beta)\gamma
 +(a\dot\tau+\dot\mu)u\Bigr]\otimes\ell.
\tag{11.19}
\]
For $(m,a)=(4,1)$, the actual invariant target generator from (11.4) is
$s(b_0\wedge b_1)\ell$.  Normalize $b_1$ so that
\[
 \gamma\wedge b_1=u_2(s)(\gamma\wedge u),
 \qquad u_2(0)=1;
\]
equivalently, use the invariant frame
\[
 \omega_2=s\,u_2(s)(\gamma\wedge u)\ell.
\tag{11.19f}
\]
The uncorrected expression $s(\gamma\wedge u)\ell$ is only invariant at
$s=0$: indeed $g_2^*(\gamma\wedge u)=\tau(z)(\gamma\wedge u)$, so its
pullback factor is $(-i)\tau(z)$.  Formula (11.11), in the genuine invariant
frame (11.19f), sends the first ordered source generator $e_0b_0$ to
minus the second coefficient of $\kappa_0(e_0)$.  Since $u_2(0)=1$,
(11.19) therefore gives the exact signed coarse value
\[
 \kappa_1(h\otimes\gamma)(p_2)
 =-\frac14\bigl(a(z_2)\dot\tau(z_2)+\dot\mu(z_2)\bigr).
\tag{11.19a}
\]

At the order-four lift $z_2$, the quotient coordinate has order four while
the modular invariant has order two at $\tau=i$; see
\cite[Chapter~VII]{Serre1973}.  Equivalently,
$j(\tau)-1728$ has a double zero at $i$ and
$j(\tau(z))-1728$ has order four at $z_2$.  Hence
\[
 \operatorname{ord}_{z_2}(\tau-i)=2,\qquad \dot\tau(z_2)=0.
\tag{11.20}
\]
By (11.7), $K(z_2)\ne0$; since
$K=6\dot\mu^2-\dot\tau\dot\beta$, (11.20) gives
$K(z_2)=6\dot\mu(z_2)^2\ne0$.  Substitution in (11.19a) therefore gives
\[
 \kappa_1(h\otimes\gamma)(p_2)=-\frac14\dot\mu(z_2)\ne0.
\tag{11.20a}
\]
The image is a global section of
$\mathcal L\otimes R^2f_*\mathcal O_X\simeq\mathcal O_B$; because it is
nonzero at $p_2$, it is a nonzero constant.  Hence
$H^0(\kappa_1)$ has rank one.
\end{proof}
